\chapter{Stato dell'arte}
\label{cap:stato-arte}

\todoscrivi{Rassegna critica, non elenco. Ogni lavoro citato deve servire a
costruire l'argomento che porta al gap di ricerca. La scheda completa di ogni
paper vive in docs/literature/: qui entra solo cio' che serve all'argomentazione.}

\section{Metodologia della revisione}
\label{sec:metodologia-revisione}
% Database interrogati, stringhe di ricerca, criteri di inclusione ed esclusione,
% intervallo temporale, numero di lavori esaminati e selezionati.
% Rende la revisione riproducibile e la difende in sede di discussione.

\section{Modelli generativi per la produzione artistica}
\label{sec:sota-modelli}

\section{Valutazione della creativita' computazionale}
\label{sec:sota-valutazione}

\section{Implicazioni etiche e autorialita'}
\label{sec:sota-etica}

\section{Lacune nella letteratura}
\label{sec:gap}
% La sezione piu' importante del capitolo: da qui deve discendere la domanda di
% ricerca dichiarata in \cref{sec:domande-ricerca}.
